\section{Experiments and analysis}
\label{sec:experiments}
The experiments below separate two questions that a single accuracy number
conflates. First, does grounding the concept space in anatomy make the
bottleneck more faithful, and what does that faithfulness cost in raw
accuracy (Tables~\ref{tab:results} and~\ref{tab:concepts})? Second, once a
model exists, does structuring the label space change \emph{how} it is wrong,
independently of \emph{how often} (Table~\ref{tab:severity-breakdown})? The
two questions have different answers: concept-space structure trades accuracy
for faithfulness and, perhaps surprisingly, does not come with safer mistakes;
label-space structure buys safer mistakes almost for free.

\begin{table}[h]
  \centering
  \footnotesize
  \setlength{\tabcolsep}{4pt}
  \caption{Top-1 accuracy and mean WordNet mistake severity
  (Eq.~\eqref{eq:sev}, lower is better) on the CUB-200-2011 test set, under
  arg\,max and CRM.}
  \label{tab:results}
  \begin{tabular}{llcccc}
    \toprule
    & & \multicolumn{2}{c}{arg\,max} & \multicolumn{2}{c}{CRM} \\
    \cmidrule(lr){3-4}\cmidrule(lr){5-6}
    Variant & Classifier & Acc. & Sev. & Acc. & Sev. \\
    \midrule
    Black-box (reference) & --- & \textbf{85.83} & \textbf{0.326} & 85.66 & 0.330 \\
    \midrule
    \multirow{2}{*}{Flat}
      & Independent & 73.35 & 0.682 & 73.32 & 0.663 \\
      & Sequential  & 79.24 & 0.531 & 78.62 & 0.544 \\
    \addlinespace
    \multirow{2}{*}{Part-conditioned}
      & Independent & 73.14 & 0.795 & 73.21 & 0.777 \\
      & Sequential  & 73.68 & 0.763 & 73.68 & 0.741 \\
    \addlinespace
    \multirow{2}{*}{\;+ hierarchical concepts}
      & Independent & 72.26 & 0.822 & 72.04 & 0.814 \\
      & Sequential  & 72.23 & 0.803 & 69.02 & 0.937 \\
    \bottomrule
  \end{tabular}
  \par\vspace{2pt}
  {\scriptsize For reference on the severity scale, the WordNet cost matrix has a
  maximum of $7.0$ and a random classifier attains an expected severity of $5.59$.}
\end{table}

\textbf{The interpretability bottleneck has a measurable cost.} The black-box
outperforms every concept-based variant on \emph{both} accuracy and severity
(Table~\ref{tab:results}), and by a wide margin---if anything an
underestimate, since its checkpoint was early-stopped at epoch $18/30$ and had
not yet converged. This is the cost of interpretability already documented by
\citet{koh2020concept}: a $109$-dimensional concept vector is a real
information bottleneck, and forcing every prediction through it necessarily
throws away whatever visual shortcuts the unconstrained network would
otherwise use, however predictive they are. No concept-based variant closes
that gap, nor should we expect one to---the bottleneck is not a free
regularizer, it is a deliberate constraint traded for an inspectable
intermediate representation. The question the rest of this section asks is
what, exactly, that trade buys.

\begin{table*}[t]
  \centering
  \footnotesize
  \setlength{\tabcolsep}{4pt}
  \caption{Mistake-severity breakdown. Sev.\ (err.) is severity computed only
  over misclassified examples; Same-genus/Same-family/Far are the share of
  errors with LCA height $1$/$2$/$\geq 3$ \citep{bertinetto2020making}.}
  \label{tab:severity-breakdown}
  \begin{tabular}{llcccccccc}
    \toprule
    & & \multicolumn{4}{c}{arg\,max} & \multicolumn{4}{c}{CRM} \\
    \cmidrule(lr){3-6}\cmidrule(lr){7-10}
    Variant & Classifier & Sev.\ (err.) & Same-genus & Same-family & Far
                         & Sev.\ (err.) & Same-genus & Same-family & Far \\
    \midrule
    Black-box (reference) & --- & \textbf{2.303} & \textbf{29.4\%} & 10.8\% & \textbf{59.8\%}
                                 & 2.300 & 29.7\% & 11.0\% & 59.3\% \\
    \midrule
    \multirow{2}{*}{Flat}
      & Independent & 2.561 & 27.6\% & 9.7\%  & 62.7\% & 2.485 & 28.2\% & 9.6\%  & 62.2\% \\
      & Sequential  & 2.557 & 25.9\% & 11.5\% & 62.6\% & 2.545 & 26.1\% & 11.4\% & 62.5\% \\
    \addlinespace
    \multirow{2}{*}{Part-conditioned}
      & Independent & 2.960 & 23.1\% & 9.8\%  & 67.1\% & 2.900 & 23.3\% & 9.9\%  & 66.8\% \\
      & Sequential  & 2.900 & 23.0\% & 11.7\% & 65.3\% & 2.815 & 24.0\% & 11.8\% & 64.2\% \\
    \addlinespace
    \multirow{2}{*}{\;+ hierarchical concepts}
      & Independent & 2.965 & 22.7\% & 9.5\%  & 67.8\% & 2.912 & 22.9\% & 9.9\%  & 67.2\% \\
      & Sequential  & 2.891 & 23.2\% & 11.6\% & 65.2\% & 3.026 & 22.6\% & 11.0\% & 66.4\% \\
    \bottomrule
  \end{tabular}
\end{table*}

\textbf{Grounding costs on both fronts at once.} Table~\ref{tab:severity-breakdown}
splits Table~\ref{tab:results}'s severity into how often each model errs and,
conditioned on erring, how bad the mistake is. A natural hope would be that
grounding trades one for the other---errs more often, perhaps, but at least
fails gracefully onto a close relative when it does. That is not what happens.
Under arg\,max, the black-box's errors are not only the least frequent but
also, on average, the mildest and the most likely to stay within the true
genus; moving to part-conditioned and hierarchical concepts makes errors both
more frequent \emph{and} proportionally farther in the taxonomy (a larger Far
share, a smaller Same-genus share) than either the black-box or the flat
model. Grounding, in other words, does not fail more gracefully---it fails
more, and slightly worse per failure. A plausible reading is that the
whole-image context grounding removes is disproportionately useful for the
hardest calls: telling apart two closely related species often comes down to
subtle, correlated cues spread across the whole bird, which is exactly the
kind of evidence a part-local crop cannot see. Confirming that mechanism,
rather than just observing its signature, is a natural next step. CRM's effect
on this breakdown is uneven across variants, which we return to below.

\textbf{The sequential classifier usually outperforms the independent one.}
Training the species classifier on the concept encoder's own predictions
(sequential) rather than on the ground-truth labels (independent) raises top-1
accuracy and lowers severity for Flat and, more modestly, for
Part-conditioned; for $+$hierarchical concepts the two regimes are
statistically indistinguishable. The explanation is straightforward: the
sequential classifier is trained on the same noisy concept probabilities it
will see at test time, so it can absorb the concept encoder's systematic
biases, whereas the independent classifier, fit on clean labels, faces a
train--test mismatch it never learns to correct. Adding the coarse-to-fine
gate evidently leaves little of that mismatch for the sequential classifier to
absorb, which is why the gap closes there. Because the sequential regime gives
each variant its best or tied-best performance, we read the effect of
grounding from that column.

\textbf{Grounding concepts reduces accuracy---by design.} Within the
sequential regime, accuracy falls step by step as we add structure: flat,
then part-conditioned, then $+$hierarchical, with severity rising in tandem.
This is expected rather than a shortcoming. The flat model's edge comes partly
from whole-image, class-correlated shortcuts that part conditioning
deliberately removes---the same shortcuts Table~\ref{tab:concepts} will show
inflating its apparent concept quality. The concept-space design trades a few
accuracy points for a bottleneck that is actually faithful: each concept is
grounded in the body region it names and cannot circularly infer the answer
from global cues, and colour/pattern concepts are additionally organised
hierarchically---an interpretability benefit no accuracy number captures.

\begin{table}[h]
  \centering
  \footnotesize
  \caption{Concept-level accuracy and macro-$F_1$ over the $109$ attributes,
  for each concept encoder alone.}
  \label{tab:concepts}
  \begin{tabular}{lccc}
    \toprule
    Concept encoder & \# concepts & Acc.\ (\%) & Macro-$F_1$ \\
    \midrule
    Flat                          & 109 & \textbf{95.66} & \textbf{0.876} \\
    Part-conditioned              & 109 & 91.36 & 0.758 \\
    \; + hierarchical concepts    & 109 & 91.24 & 0.753 \\
    \bottomrule
  \end{tabular}
\end{table}

\textbf{Higher concept scores do not imply more faithful concepts.}
Table~\ref{tab:concepts} evaluates concept quality in isolation, and the flat
encoder comes out well ahead of the part-conditioned one---even though the
latter is the one that actually breaks prediction circularity. Far from
contradicting our thesis, this gap \emph{supports} it: the flat encoder can
infer an attribute from whole-image cues that are diagnostic of the species
(e.g.\ recognising \emph{wing color} from the overall plumage pattern) rather
than from the region the attribute names, the same shortcut mechanism that
inflates its species accuracy in Table~\ref{tab:results}. The part-conditioned
encoder, forced to work from the crop of the relevant body part alone, cannot
take that shortcut; its lower score is therefore a more honest measurement of
whether the concept is grounded in its own evidence, not a worse one. A
high concept score, on its own, is consequently not evidence of a faithful
bottleneck---it can equally be evidence of leakage. Adding the coarse-to-fine
gate barely moves the number either way: the extra structure is essentially
free.

\textbf{CRM buys lower severity almost everywhere, for almost no accuracy.}
Applied exactly as \citet{karthik2021nocost} define it---no retraining,
no extra parameters, no free choices---CRM lowers severity for five of the six
concept-based rows in Table~\ref{tab:results}, in every case at a cost of well
under half an accuracy point. It leaves the black-box and Flat-Sequential
essentially unchanged: both already have sharp, confident posteriors, so there
is little uncertainty left for a severity-aware re-ranking to act on.
$+$Hierarchical-Sequential is the one row where CRM makes both accuracy and
severity worse. A re-ranking rule built from expected severity is only as
good as the posterior it re-ranks; when that posterior is not a faithful
estimate of the model's own uncertainty, minimising \emph{expected} severity
under it can easily miss actual severity. CRM's benefit is therefore real but
conditional on the classifier it is layered onto, not a property of the
decision rule alone---exactly the caveat a severity-aware method should be
reported with.
